\documentclass[11pt]{article}

\usepackage[a4paper,margin=1in]{geometry}
\usepackage{amsmath,amssymb,bm}
\usepackage{booktabs}
\usepackage{hyperref}

\hypersetup{
    colorlinks=true,
    linkcolor=blue,
    citecolor=blue,
    urlcolor=blue
}

\newcommand{\master}{geqoe\_averaged\_j2\_zonal\_note}

\title{Taylor Integration of the GEqOE Mean-Element ODE\\
\large Heyoka-Based Propagation of the Averaged Zonal Model}
\author{Jose Manuel Montilla Garc\'ia}
\date{\today}

\begin{document}
\maketitle

\begin{abstract}
This companion note describes three strategies for accelerating the
mean-element propagation step of the GEqOE averaged zonal theory (master
note, \texttt{\master{}}). The mean dynamics for the general mixed-zonal
model ($J_2$--$J_5$) constitute a smooth, slowly varying, low-dimensional
autonomous ODE whose right-hand side is a known closed-form function. This
ODE is an ideal candidate for high-order Taylor integration via the heyoka
library, which provides automatic differentiation of the Taylor coefficients
and adaptive step-size control. Three progressively more ambitious paths are
described: (A)~direct Taylor integration of the 4D mean ODE,
(B)~JIT-compiled evaluation of the mean drift right-hand side via heyoka
\texttt{cfunc}, and (C)~parametric Taylor maps for reusable propagation
across orbit families.
\end{abstract}

\tableofcontents

\section{The Mean-Element ODE}

\subsection{Structure}

The master note (Section~10.5) establishes that the total-order-$N$ averaged
slow system, after summation over all zonal degrees $n = 2, \ldots, N$, is
\begin{align}
\dot{\bar g}
&=
\sum_{\substack{1 \le m \le N \\ m\;\mathrm{odd}}}
G^{(c)}_m(\bar\nu, \bar g, \bar Q)\cos(m\bar\omega)
+
\sum_{\substack{2 \le m \le N \\ m\;\mathrm{even}}}
G^{(s)}_m(\bar\nu, \bar g, \bar Q)\sin(m\bar\omega),
\label{eq:gdot}\\
\dot{\bar Q}
&=
\sum_{\substack{1 \le m \le N \\ m\;\mathrm{odd}}}
Q^{(c)}_m(\bar\nu, \bar g, \bar Q)\cos(m\bar\omega)
+
\sum_{\substack{2 \le m \le N \\ m\;\mathrm{even}}}
Q^{(s)}_m(\bar\nu, \bar g, \bar Q)\sin(m\bar\omega),
\label{eq:Qdot}\\
\dot{\bar\Psi}
&=
P_0(\bar\nu, \bar g, \bar Q)
+
\sum_{\substack{1 \le m \le N \\ m\;\mathrm{odd}}}
P^{(s)}_m \sin(m\bar\omega)
+
\sum_{\substack{2 \le m \le N \\ m\;\mathrm{even}}}
P^{(c)}_m \cos(m\bar\omega),
\label{eq:Psidot}\\
\dot{\bar\Omega}
&=
N_0(\bar\nu, \bar g, \bar Q)
+
\sum_{\substack{1 \le m \le N \\ m\;\mathrm{odd}}}
N^{(s)}_m \sin(m\bar\omega)
+
\sum_{\substack{2 \le m \le N \\ m\;\mathrm{even}}}
N^{(c)}_m \cos(m\bar\omega),
\label{eq:Omegadot}
\end{align}
where $\bar\omega = \bar\Psi - \bar\Omega$ is the mean argument of
pericenter. The coefficient functions $G^{(c)}_m$, $P_0$, etc., are known
closed-form algebraic expressions in $(\bar\nu, \bar g, \bar Q)$ involving:
\begin{itemize}
\item rational functions of $\bar g$ and $\bar Q$ (with denominators
involving powers of $1 - \bar g^2$ and $1 + \bar Q^2$);
\item the auxiliary $\bar\beta = \sqrt{1 - \bar g^2}$, which introduces a
single square root;
\item the semi-major axis $\bar a = (\mu/\bar\nu^2)^{1/3}$.
\end{itemize}

\subsection{Properties relevant to Taylor integration}

\begin{enumerate}
\item \textbf{Low dimension}: 4 state variables
$(\bar g, \bar Q, \bar\Psi, \bar\Omega)$, with $\bar\nu$ constant.

\item \textbf{Smooth}: the right-hand side is $C^\infty$ (in fact, real
analytic) for $0 < \bar g < 1$ and $\bar Q > 0$. There are no
discontinuities, switches, or stiff terms.

\item \textbf{Slowly varying}: the characteristic timescale is the
pericenter precession period, which for typical LEO/MEO orbits is hundreds
to thousands of orbital periods.

\item \textbf{Autonomous}: no explicit time dependence (a direct consequence
of the zonal-only restriction).

\item \textbf{Finite trigonometric}: the dependence on $\bar\omega$ is a
finite Fourier series with at most $N - 2 = 3$ harmonics (for $N = 5$).
\end{enumerate}

These properties make the mean ODE an ideal target for high-order Taylor
integration: the smoothness guarantees rapid convergence of the Taylor
series, the low dimension keeps the cost per step negligible, and the slow
variation allows enormous time steps.

\subsection{The fast phase}

In addition to the slow state, the mean fast phase
$\bar M = \bar L - \bar\Psi$ must be advanced. Its rate is
\begin{equation}
\dot{\bar M}
= \bar\nu
+ \dot{\bar L}_{\mathrm{sec}} - \dot{\bar\Psi},
\end{equation}
where $\dot{\bar L}_{\mathrm{sec}}$ includes the mean-motion correction
from the zonal perturbation. Because $\bar\nu$ is constant and
$\dot{\bar\Psi}$ is already part of the slow ODE, $\bar M$ can be integrated
as a 5th state variable or computed by quadrature after the slow propagation.

For the pure $J_2$ case, $\dot{\bar M}$ is constant and the integration is
trivial: $\bar M(t) = \bar M_0 + \dot{\bar M} \cdot (t - t_0)$. For the
general zonal case, $\dot{\bar M}$ depends on $\bar g$ and $\bar Q$ (which
vary slowly), so it must be integrated alongside the slow state.


\section{Path A: Direct Taylor Integration}

\subsection{Concept}

Build the mean ODE right-hand side \eqref{eq:gdot}--\eqref{eq:Omegadot} as
a heyoka expression tree and integrate it with heyoka's adaptive
\texttt{taylor\_adaptive} integrator.

\subsection{Implementation}

\begin{enumerate}
\item \textbf{Express the coefficient functions as heyoka expressions.}
The symbolic expressions for $G^{(c)}_m(\bar\nu, \bar g, \bar Q)$, etc.,
are currently stored as SymPy expressions. These can be converted to heyoka
expression trees either by:
\begin{itemize}
\item direct translation (mapping SymPy operators to heyoka operators);
\item or by generating Python code that builds the heyoka expression tree
from the SymPy output.
\end{itemize}
The key operations are: arithmetic $(+, -, \times, /)$, powers
(integer and half-integer via $\sqrt{\cdot}$), and trigonometric functions
$(\sin, \cos)$ of $\bar\omega = \bar\Psi - \bar\Omega$.

\item \textbf{Define the state variables.} In heyoka:
\begin{verbatim}
g, Q, Psi, Omega, M = hy.make_vars("g", "Q", "Psi", "Omega", "M")
omega = Psi - Omega
\end{verbatim}
The mean frequency $\bar\nu$ enters as a parameter
(\texttt{hy.par[0]}).

\item \textbf{Build the ODE system.}
\begin{verbatim}
sys = [
    (g,     g_dot_expr),
    (Q,     Q_dot_expr),
    (Psi,   Psi_dot_expr),
    (Omega, Omega_dot_expr),
    (M,     M_dot_expr),
]
ta = hy.taylor_adaptive(sys, ic, pars=[nu_val])
\end{verbatim}

\item \textbf{Propagate.} Call \texttt{ta.propagate\_until(t\_final)} or
\texttt{ta.propagate\_grid(t\_grid)}. The adaptive Taylor integrator will
automatically select step sizes appropriate to the slow dynamics---likely
spanning tens to hundreds of orbital periods per step.
\end{enumerate}

\subsection{Expected performance}

For a Taylor order of 15--20 (heyoka's default adaptive range), the Taylor
series converges to machine precision in a single evaluation when the step
size is well within the convergence radius. For the mean ODE, the
convergence radius is set by the slowest precession timescale, which is
vastly longer than the orbital period.

\textbf{Rough estimate}: For a LEO orbit with $J_2$ pericenter precession
period $\sim$60 days, the Taylor integrator might take $\sim$10 steps to
cover one precession cycle, compared to $\sim$1000 orbital periods. Each
step involves evaluating $\sim$20 Taylor coefficients for 5 state variables,
at a cost of $\sim$1000 floating-point operations. Total cost for one
precession cycle: $\sim 10^4$ FLOPs---essentially instantaneous.

\subsection{Variational equations}

For covariance propagation or sensitivity analysis, heyoka's
\texttt{var\_ode\_sys} can automatically generate the variational equations
of the mean ODE:
\begin{verbatim}
vsys = hy.var_ode_sys(sys, hy.var_args.vars, order=1)
\end{verbatim}
This gives the $5 \times 5$ state transition matrix of the mean dynamics at
no additional implementation cost. The mean STM would capture the
sensitivity of the mean state to initial conditions, which is useful for
orbit determination in mean elements.


\section{Path B: Compiled Function Evaluation}

\subsection{Concept}

Even without Taylor integration, the mean drift right-hand side can be
compiled into a fast, vectorized evaluation function via heyoka's
\texttt{cfunc} (compiled function) facility. This would replace the current
Python-level evaluation of the symbolic expressions with JIT-compiled machine
code.

\subsection{Implementation}

\begin{enumerate}
\item Build the heyoka expression tree for each component of the mean drift:
$\dot{\bar g}$, $\dot{\bar Q}$, $\dot{\bar\Psi}$, $\dot{\bar\Omega}$,
$\dot{\bar M}$.

\item Create a compiled function:
\begin{verbatim}
cf = hy.cfunc(
    [g_dot_expr, Q_dot_expr, Psi_dot_expr,
     Omega_dot_expr, M_dot_expr],
    vars=[g, Q, Psi, Omega, M],
    pars=[hy.par[0]],  # nu
)
\end{verbatim}

\item Evaluate at any state:
\begin{verbatim}
rhs = cf([g_val, Q_val, Psi_val, Omega_val, M_val],
         pars=[nu_val])
\end{verbatim}
\end{enumerate}

\subsection{Use cases}

\begin{itemize}
\item \textbf{Classical integrator backend}: Use the compiled RHS with
\texttt{scipy.integrate.solve\_ivp} or a custom Runge--Kutta integrator.
This avoids the overhead of Python-level symbolic evaluation while keeping
the integrator choice flexible.

\item \textbf{Batch evaluation}: \texttt{cfunc} supports SIMD-vectorized
evaluation over arrays of states, enabling efficient Monte Carlo or ensemble
propagation.

\item \textbf{Gradient computation}: heyoka's compiled functions support
automatic differentiation, which can provide Jacobians for
Newton-type solvers (e.g., orbit determination, targeting).
\end{itemize}

\subsection{Expected speedup}

The current mean drift evaluation is implemented in Python using cached
SymPy \texttt{lambdify} functions. A heyoka \texttt{cfunc} would replace
this with compiled C code, providing a speedup of $\sim$10--100$\times$ per
evaluation. For batch propagation of thousands of orbits, the SIMD
vectorization would add another $\sim$4--8$\times$ factor.

\subsection{Short-period kernel compilation}

The same approach can be applied to the short-period correction functions
$\mathcal{G}_1(G)$, $\mathcal{Q}_1(G)$, $\mathcal{P}_1(G)$,
$\mathcal{N}_1(G)$, $\ell_1(K)$. These are rational functions of
$\sin G$, $\cos G$ (or $\sin K$, $\cos K$) with coefficients in $(q, Q)$.
Compiling them via \texttt{cfunc} would make the mean-to-osculating
reconstruction step equally fast.


\section{Path C: Parametric Taylor Maps}

\subsection{Concept}

For applications requiring propagation of many orbits with similar
characteristics (e.g., constellation design, catalog maintenance), it may be
advantageous to precompute Taylor maps that are valid over a neighborhood of
the initial-condition space, rather than integrating each orbit individually.

\subsection{Background: Taylor maps}

A Taylor map of order $p$ around a reference state $\bm x_0$ at time $t_0$
is a polynomial approximation to the flow:
\begin{equation}
\bm x(t_0 + \Delta t)
= \sum_{|\bm k| = 0}^{p}
\bm c_{\bm k}(\Delta t)\,
(\bm x_0 - \bm x_{\mathrm{ref}})^{\bm k},
\end{equation}
where $\bm k$ is a multi-index and $\bm c_{\bm k}(\Delta t)$ are the Taylor
coefficients in the initial-condition deviation. These coefficients can be
computed by propagating the variational equations to order $p$.

\subsection{Application to the mean ODE}

The mean ODE has only 4 slow state variables (with $\bar\nu$ as a parameter).
A Taylor map of order $p = 4$ in the initial conditions would involve
$\binom{4 + 4}{4} = 70$ coefficients per state variable, or 350 total. This
is a modest number.

For a given reference orbit class (e.g., sun-synchronous at 700~km altitude),
the Taylor map would provide the mean-state evolution for any orbit in the
neighborhood without re-integration. The validity domain depends on how much
the dynamics vary across the orbit family; for orbits with similar $a$, $e$,
$i$, the Taylor map should be valid over the range of initial $\omega$ values.

\subsection{Heyoka implementation}

Heyoka's \texttt{var\_ode\_sys} with \texttt{order=p} automatically generates
the variational equations to order $p$. For $p = 4$ and 5 state variables,
this produces $5 + 5 \times 5 + 5 \times 15 + 5 \times 35 + 5 \times 70
= 630$ ODEs. While larger than the original 5-variable system, this is still
a small ODE by heyoka standards (heyoka routinely handles thousands of
variables).

\begin{verbatim}
vsys = hy.var_ode_sys(sys, hy.var_args.vars, order=4)
ta = hy.taylor_adaptive(vsys, ic_extended)
ta.propagate_until(t_final)
# Extract Taylor map coefficients from ta.state
\end{verbatim}

\subsection{Caching and reuse}

The Taylor map coefficients could be:
\begin{itemize}
\item \textbf{Precomputed and stored} for a grid of reference orbits
(parameterized by $a$, $e$, $i$). Each grid point stores 350 floating-point
numbers per time step.

\item \textbf{Evaluated by polynomial Horner evaluation} for any orbit in
the neighborhood. This is $O(p^d)$ FLOPs with $d = 4$ dimensions and
$p = 4$, i.e., $\sim$256 FLOPs per evaluation---negligible.

\item \textbf{Composed for multi-step propagation}: Taylor maps from
consecutive time intervals can be composed (polynomial composition) to extend
the propagation window without re-integration.
\end{itemize}

\subsection{Limitations}

\begin{enumerate}
\item The Taylor map is only accurate within its convergence domain. For
orbits that differ significantly from the reference (e.g., large eccentricity
difference), the map must be recomputed.

\item The variational equations at high order grow rapidly. Order $p = 4$
with 5 variables gives 630 ODEs, which is manageable; order $p = 6$ would
give $\sim$2500 ODEs, which is still feasible but slower to compute.

\item The mean ODE is only 4-dimensional, so the advantage of Taylor maps
over direct integration is primarily in the amortization of the integration
cost over many orbits, not in the per-orbit cost (which is already
negligible with Path A).
\end{enumerate}


\section{Comparison of Paths}

\begin{center}
\begin{tabular}{lccc}
\toprule
& Path A & Path B & Path C \\
& Taylor ODE & Compiled RHS & Taylor maps \\
\midrule
Implementation effort & Low & Low & Medium \\
Per-orbit propagation cost & Negligible & Low & Negligible \\
Batch/ensemble suitability & Good & Excellent & Excellent \\
Covariance/STM & Automatic & Manual & Built-in \\
Reusability across orbits & None & None & High \\
Prerequisite & heyoka & heyoka & heyoka + grid \\
\bottomrule
\end{tabular}
\end{center}

\textbf{Recommendation}: Implement Path A first (lowest effort, highest
immediate payoff). Path B is a natural companion for the short-period kernel
evaluation. Path C is a longer-term investment for catalog-scale
applications.


\section{Implementation Roadmap}

\begin{enumerate}
\item \textbf{SymPy-to-heyoka translator}: Write a utility that converts the
stored SymPy coefficient functions ($G^{(c)}_m$, $P_0$, etc.) into heyoka
expression trees. This is a one-time infrastructure investment that also
benefits the short-period kernel compilation.

\item \textbf{Path A prototype}: Build the 5-variable mean ODE in heyoka
and validate against the current Python-based mean propagator. Verify that
the Taylor step sizes are as large as expected ($\gg$ orbital period).

\item \textbf{Path B prototype}: Compile the mean drift RHS and the
short-period kernels as \texttt{cfunc}. Benchmark against the current
\texttt{lambdify}-based evaluation.

\item \textbf{Path A + variational}: Add \texttt{var\_ode\_sys} to the
Path A integrator and validate the mean STM against finite differences.

\item \textbf{Path C exploration}: For a representative orbit family (e.g.,
sun-synchronous LEO), compute a 4th-order Taylor map and assess the validity
domain and accuracy.
\end{enumerate}

\end{document}
